\documentclass[11pt,a4paper]{article}

% -----------------------------
% Kodierung & Sprache
% -----------------------------
\usepackage[T1]{fontenc}
\usepackage[utf8]{inputenc}
\usepackage[ngerman]{babel}
\usepackage{lmodern}
\usepackage{microtype}

% -----------------------------
% Mathematik
% -----------------------------
\usepackage{amsmath,amssymb,amsfonts,bm,mathtools}

% -----------------------------
% Layout, Grafik, Links
% -----------------------------
\usepackage[a4paper,margin=2.5cm]{geometry}
\usepackage{graphicx}
\usepackage{xcolor}
\usepackage{hyperref}
\usepackage{float}
\usepackage{booktabs}
\usepackage{caption}

\hypersetup{
	colorlinks=true,
	linkcolor=blue,
	citecolor=blue,
	urlcolor=blue,
	pdftitle={Ein quaternionischer Dirac-Operator auf Primzahlvierlingen},
	pdfauthor={Thomas Hoffbauer}
}

\title{Ein quaternionischer Dirac-Operator auf Primzahlvierlingen\\
	und seine spektrale Beziehung zu den Riemannschen Nullstellen}

\author{Thomas Hoffbauer}
\date{\today}

\begin{document}
	
	\maketitle
	
	\begin{abstract}
		Wir konstruieren einen diskreten Operator auf der Menge der Primzahlvierlinge
		\[
		(p,\; p+2,\; p+6,\; p+8),
		\]
		indem wir diese in eine quaternionische Struktur einbetten.
		Übergänge zwischen aufeinanderfolgenden Vierlingen definieren eine
		nichtkommutative Dynamik, aus der ein hermitescher Operator gewonnen wird.
		
		Numerische Experimente zeigen eine starke Korrelation zwischen dem resultierenden
		Spektrum und den nichttrivialen Nullstellen der Riemannschen Zetafunktion.
		Darüber hinaus weist die Phasenstruktur der Übergänge eine Clusterbildung auf,
		die mit einer SU(3)-artigen Symmetrie konsistent ist und auf eine intrinsische
		algebraische Organisation der Primzahlen hindeutet.
	\end{abstract}
	
	\section{Einleitung}
	
	Die Verteilung der Primzahlen gehört zu den zentralen Problemen der Zahlentheorie.
	Die Hilbert--Pólya-Vermutung legt nahe, dass die nichttrivialen Nullstellen der
	Riemannschen Zetafunktion als Spektrum eines geeigneten selbstadjungierten Operators
	interpretiert werden können.
	
	In dieser Arbeit konstruieren wir einen solchen Operator auf Primzahlvierlingen
	unter Verwendung quaternionischer Geometrie.
	
	\section{Primzahlvierlinge}
	
	Wir betrachten Primzahlvierlinge der Form
	\[
	(p,\; p+2,\; p+6,\; p+8).
	\]
	
	Jeder Vierling wird in ein Quaternion eingebettet:
	\[
	Q = (p,\; p+2,\; p+6,\; p+8).
	\]
	
	\section{Quaternionischer Übergangsoperator}
	
	Für zwei aufeinanderfolgende Vierlinge \(Q_i\) und \(Q_j\) definieren wir
	\[
	\Delta_{ij} = Q_i^{-1} Q_j.
	\]
	
	Die relevanten Observablen sind
	\[
	\|\Delta_{ij}\|
	\qquad \text{und} \qquad
	\theta_{ij}
	=
	\arccos\!\left(
	\frac{\mathrm{Re}(\Delta_{ij})}{\|\Delta_{ij}\|}
	\right).
	\]
	
	Aus diesen Größen definieren wir zunächst das komplexe Übergangsgewicht
	\[
	w_{ij}
	=
	\frac{1}{\|\Delta_{ij}\|}\, e^{i\theta_{ij}}.
	\]
	
	Ein hermitescher Operator ergibt sich dann durch
	\[
	D_{ij} = w_{ij}, \qquad D_{ji} = \overline{w_{ij}}.
	\]
	
	Damit ist \(D\) per Konstruktion hermitesch.
	
	\section{Numerische Ergebnisse}
	
	Wir berechnen die Eigenwerte von \(D\) und vergleichen diese mit den
	Imaginärteilen der nichttrivialen Nullstellen der Riemannschen Zetafunktion.
	
	\subsection{Korrelation}
	
	Die beobachtete Korrelation beträgt näherungsweise
	\[
	\mathrm{corr}(\lambda_n,\gamma_n) \approx 0.97.
	\]
	
	\subsection{Abstandsstatistik}
	
	Die Verteilung der nächsten Nachbarabstände wird mit den Vorhersagen
	des Gaussian Unitary Ensemble (GUE) verglichen:
	\[
	P(s) \sim s^2 e^{-c s^2}.
	\]
	
	\subsection{Paar-Korrelation}
	
	Wir vergleichen mit der Montgomeryschen Paar-Korrelationsfunktion:
	\[
	R_2(u)
	=
	1-
	\left(
	\frac{\sin(\pi u)}{\pi u}
	\right)^2.
	\]
	
	\section{SU(3)-artige Struktur}
	
	Primzahlen werden nach Restklassen modulo \(12\) klassifiziert:
	\[
	\{1,\,5,\,7,\,11\}.
	\]
	
	Übergänge zwischen diesen Klassen zeigen eine strukturierte Phasenverteilung
	mit Clustern nahe Winkeln von \(2\pi/3\), was auf eine zyklische Symmetrie
	analog zu SU(3) hindeutet.
	
	\section{Abbildungen}
	
	Abbildung~\ref{fig:theta} zeigt die Verteilung der Übergangsphasen.
	Abbildung~\ref{fig:spacing} zeigt die normierte Abstandsstatistik.
	Abbildung~\ref{fig:paircorr} vergleicht die empirische Paar-Korrelation
	mit der Montgomery-Kurve.
	
	\begin{figure}[H]
		\centering
		\includegraphics[width=0.78\textwidth]{theta_hist.png}
		\caption{Histogramm der Übergangsphasen \(\theta_{ij}\). Peaks nahe
			\(2\pi/3\) wären mit einer SU(3)-artigen Zyklusstruktur vereinbar.}
		\label{fig:theta}
	\end{figure}
	
	\begin{figure}[H]
		\centering
		\includegraphics[width=0.78\textwidth]{spacing_hist.png}
		\caption{Normierte Nearest-Neighbor-Abstände des Spektrums im Vergleich
			zur GUE-Erwartung.}
		\label{fig:spacing}
	\end{figure}
	
	\begin{figure}[H]
		\centering
		\includegraphics[width=0.78\textwidth]{pair_correlation.png}
		\caption{Empirische Paar-Korrelationsfunktion des Spektrums im Vergleich
			zur Montgomery-Formel.}
		\label{fig:paircorr}
	\end{figure}
	
	\section{Diskussion}
	
	Die Konstruktion liefert einen Kandidaten für einen diskreten Dirac-Operator
	auf Primzahlkonfigurationen.
	
	Die spektrale Übereinstimmung mit den Riemannschen Nullstellen sowie das Auftreten
	interner Symmetriestrukturen deuten auf eine tiefere geometrische Organisation
	der Primzahlen hin. Zugleich muss betont werden, dass hohe numerische Korrelationen
	allein noch keinen strengen mathematischen Zusammenhang beweisen; entscheidend sind
	Robustheitstests, Vergleichsmodelle und die Stabilität gegenüber Modellvarianten.
	
	\section{Schlussfolgerung}
	
	Wir haben einen quaternionischen Operator auf Primzahlvierlingen eingeführt,
	dessen Spektrum statistische und numerische Übereinstimmungen mit den
	Riemannschen Nullstellen zeigt.
	
	Weitere Arbeiten sind erforderlich, um die analytischen Eigenschaften dieses
	Operators rigoros zu untersuchen und seine Beziehung zu etablierten
	zahlentheoretischen Rahmenwerken zu klären.
	
\end{document}